\chapter[Local Diffeomorphism]{Local Diffeomorphism}\label{ch:sec:local-diffeomorphism}
Throughout this chapter let
\[
F = Q \circ L : \Delta_k^\circ \longrightarrow \operatorname{Im}(Q)
\]
denote the canonical logarithmic coordinate map, whose components were introduced in Chapter~\ref{ch:sec:the-posterior-manifold} and which is named and studied in Chapter~\ref{ch:sec:canonical-projection-and-differential-factorisation}.
The objective of this chapter is to prove that \(F\) is a local diffeomorphism.

\begin{proposition}[Smoothness]
\label{ch:prop:5-1}
The map \(F\) is smooth.
\end{proposition}
\begin{proof}
The map \(L\) is smooth because each component is a logarithm of a strictly positive coordinate.
The projection \(Q\) is linear, hence smooth.
The composition \(F = Q \circ L\) is therefore smooth. 
\end{proof}

This chapter establishes that the logarithmic coordinates are a valid chart on the posterior manifold. This justifies every local computation---gradients, Hessians, Taylor expansions---that appears in the stability analysis (Chapter~\ref{ch:sec:stability-of-the-universal-optimality-defect}) and the variational theory (Chapters~20--25).

\section{Differential}\label{ch:sec:differential}
\begin{proposition}[Differential of the Logarithmic Map]
\label{ch:prop:5-2}
For every \(P \in \Delta_k^\circ\),
$DF_P = Q \circ DL_P,$
where \(DL_P(v) = (v_i / P_i)_{i=1}^k\).
\end{proposition}
\begin{proof}
By the chain rule,
$DF_P(v) = Q(DL_P(v)).$
The expression for \(DL_P\) follows from the componentwise derivative of the logarithm. 
\end{proof}
\begin{corollary}[Invertibility on the Tangent Space]
\label{ch:cor:5-3}
The differential \(DF_P\) is a linear isomorphism between \(T_P \Delta_k^\circ\) and \(\operatorname{Im}(Q)\).
\end{corollary}
\begin{proof}
By Proposition~\ref{ch:prop:5-2}, \(DF_P = Q \circ DL_P\).
By Corollary~\ref{ch:cor:3-3}, \(DL_P : T_P \Delta_k^\circ \to W_P\) is a linear isomorphism.
By Corollary~\ref{ch:cor:4-3}, the restriction \(Q|_{W_P} : W_P \to \operatorname{Im}(Q)\) is a linear isomorphism.
Since the composition of linear isomorphisms is again a linear isomorphism,
\(DF_P\) is itself a linear isomorphism. 
\end{proof}
\section{Local Diffeomorphism Property}\label{ch:sec:local-diffeomorphism-property}
\begin{theorem}[Local Diffeomorphism]
\label{ch:thm:5-4}
The canonical logarithmic coordinate map
\[
F : \Delta_k^\circ \longrightarrow \operatorname{Im}(Q)
\]
is a local diffeomorphism.
\end{theorem}
\begin{proof}
By Proposition~\ref{ch:prop:5-1}, the map \(F\) is smooth.
By Proposition~\ref{ch:prop:5-2}, the differential \(DF_P\) is invertible at every point \(P \in \Delta_k^\circ\).
By the Inverse Function Theorem, \(F\) is a local diffeomorphism. 
\end{proof}
\section{Consequences}\label{ch:sec:consequences-diffeomorphism_local}
\begin{corollary}[Local Logarithmic Coordinates]
\label{ch:cor:5-5}
The logarithmic coordinate map provides a local chart on the posterior manifold. Consequently, every smooth function on the posterior manifold can be expressed locally in logarithmic coordinates.
\end{corollary}
\begin{proof}
Since \(F\) is a local diffeomorphism, its inverse restricted to a neighborhood of each point provides a local coordinate chart.
The second statement follows from the standard fact that local diffeomorphisms allow pulling back smooth functions. 
\end{proof}

\paragraph{Running example: geo-indistinguishability.}
The local diffeomorphism theorem guarantees that for geo-indistinguishability (Section~\ref{ch:sec:geo-indistinguishability-and-location-privacy}) the logarithmic quotient coordinates provide a valid local chart around every posterior.  In a neighbourhood of any true location $P$, the canonical map $F=Q\circ L$ has differential $DF_P(v)=Q(v_i/P_i)$ applied componentwise, and the UOD is $\mathcal D(P,Q)=\|F(P)-F(Q)\|^2$.  For small perturbations $\delta=P-Q$, the first-order term of $F(P)-F(Q)$ is $Q(\delta_i/P_i)$, so $\mathcal D\approx\sum_i \delta_i^2/P_i$ up to the quotient projection---a weighted Euclidean distance in which small probabilities amplify the cost of perturbing them.

This result establishes the logarithmic quotient coordinates as a valid local coordinate system for differential geometry on the posterior manifold. The next chapter introduces the pullback metric induced by this local diffeomorphism and develops the associated Riemannian geometry.
